\documentclass[12pt,a4paper,fleqn]{article}

%------------------------------------------------
% Кодировки и шрифты
%------------------------------------------------
\usepackage[utf8]{inputenc}
\usepackage[T2A]{fontenc}
\usepackage[russian]{babel}
\usepackage{amsmath,amssymb,mathtools}

%------------------------------------------------
% Оформление, отступы и выравнивание
%------------------------------------------------
\usepackage{geometry}
\geometry{left=20mm,right=20mm,top=20mm,bottom=20mm} 

\usepackage{setspace}
\onehalfspacing 

% Пакет для настройки списков
\usepackage{enumitem} 
\setlist[enumerate,itemize]{leftmargin=*} % Убирает отступы у списков

\raggedright % Прижимает весь текст влево (без выравнивания по ширине)
\setlength{\parindent}{0pt} % Убирает красную строку
\setlength{\parskip}{1em} % Отступ между абзацами
\setlength{\mathindent}{0pt} % Прижимает формулы влево

\allowdisplaybreaks 

\begin{document}

\noindent\textbf{\Large Задача 10 (Часть 2) — Вариант 17}\\[0.5cm]

\section*{1. Постановка задачи}

Найти экстремум функционала:
\begin{flalign*}
& J[x]=\int_0^\pi(\dot x^2-x^2)\,dt\to\operatorname{extr}, &
\end{flalign*}
при условиях:
\begin{flalign*}
& x(0)=0,\qquad \int_0^\pi x\cos t\,dt=\frac\pi2,\qquad \int_0^\pi x\sin t\,dt=-2. &
\end{flalign*}

\section*{2. Сведение к задаче Лагранжа}

Введем замену переменной \(\dot x=u\). Тогда функционал примет вид:
\begin{flalign*}
& \int_0^\pi(u^2-x^2)\,dt\to\operatorname{extr}, &
\end{flalign*}
с дополнительной связью и граничными условиями:
\begin{flalign*}
& \dot x-u=0,\quad x(0)=0,\quad \int_0^\pi x\cos t\,dt=\frac\pi2,\quad \int_0^\pi x\sin t\,dt=-2. &
\end{flalign*}

Функция Лагранжа:
\begin{flalign*}
& L=\lambda_0(u^2-x^2)+p(t)(\dot x-u)-\lambda_1 x\cos t-\lambda_2 x\sin t, & \\
& l=\lambda_1\frac\pi2-2\lambda_2+\lambda_3 x(0). &
\end{flalign*}
Множители: \(\lambda_0,\lambda_1,\lambda_2,\lambda_3\) и \(p(t)\).

\section*{3. Необходимые условия}

\begin{enumerate}
    \item \textbf{Стационарность по \(x\):}
    \begin{flalign*}
    & \frac{d}{dt}L_{\dot x}=L_x \implies \dot p=-2\lambda_0 x-\lambda_1\cos t-\lambda_2\sin t. &
    \end{flalign*}
    
    \item \textbf{Стационарность по \(u\):}
    \begin{flalign*}
    & L_u=0 \implies p=2\lambda_0 u. &
    \end{flalign*}
    
    \item \textbf{Условия трансверсальности:}
    \begin{flalign*}
    & p(0)=\lambda_3,\qquad p(\pi)=0. &
    \end{flalign*}
    
    \item \textbf{Согласование знаков:} \(\lambda_0\ge0\); \(\lambda_1,\lambda_2,\lambda_3\) любого знака.
\end{enumerate}

\section*{4. Исследование нормальности}

\noindent\textbf{Случай \(\lambda_0=0\):}
\begin{flalign*}
& p \equiv 0 \implies \lambda_1\cos t+\lambda_2\sin t\equiv0 \implies \lambda_1=\lambda_2=0 & \\
& \implies \lambda_3=0 \implies \text{все множители нулевые — невозможно.} &
\end{flalign*}

\noindent\textbf{Случай \(\lambda_0\neq 0\):}
\begin{flalign*}
& \lambda_0=1 \implies p=2\dot x, \quad \ddot x+x=-\frac{\lambda_1}{2}\cos t-\frac{\lambda_2}{2}\sin t, & \\
& x(0)=0, \quad \dot x(\pi)=0. &
\end{flalign*}

\section*{5. Решение краевой задачи}

\begin{flalign*}
& \ddot x + x = -\frac{\lambda_1}{2}\cos t - \frac{\lambda_2}{2}\sin t \equiv 2a\cos t - 2b\sin t, & \\
& a = -\frac{\lambda_1}{4}, \quad b = \frac{\lambda_2}{4} & \\
& x(t) = C_1\cos t + C_2\sin t + at\sin t + bt\cos t & \\
& x(0) = 0 \implies C_1 = 0 & \\
& \dot x(t) = C_2\cos t + a(\sin t + t\cos t) + b(\cos t - t\sin t) & \\
& \dot x(\pi) = 0 \implies -C_2 - \pi a - b = 0 \implies C_2 = -\pi a - b &
\end{flalign*}

\section*{6. Интегральные условия}

Вычислим необходимые интегралы:
\begin{flalign*}
& \int_0^\pi\sin t\cos t\,dt=0,\qquad \int_0^\pi t\sin t\cos t\,dt=-\frac\pi4, & \\
& \int_0^\pi\sin^2 t\,dt=\frac\pi2,\qquad \int_0^\pi t\cos^2 t\,dt=\int_0^\pi t\sin^2 t\,dt=\frac{\pi^2}4. &
\end{flalign*}

Подставляем \(x(t)\) в заданные интегральные ограничения:
\begin{flalign*}
& \int_0^\pi x\cos t\,dt=-\frac\pi4 a+\frac{\pi^2}4 b=\frac\pi2 \quad\Rightarrow\quad -a+\pi b=2. & \\
& \int_0^\pi x\sin t\,dt=\frac\pi2 C_2+\frac{\pi^2}4 a-\frac\pi4 b=-2. &
\end{flalign*}

Подставим \(C_2=-\pi a-b\) во второе уравнение:
\begin{flalign*}
& \pi^2 a+3\pi b=8. &
\end{flalign*}
Из первого уравнения \(a=\pi b-2\). Подставляем:
\begin{flalign*}
& \pi(\pi^2+3)b=2(\pi^2+4) \implies b=\frac{2(\pi^2+4)}{\pi(\pi^2+3)},\qquad a=\frac2{\pi^2+3}. &
\end{flalign*}
Следовательно:
\begin{flalign*}
& C_2=-\frac{4(\pi^2+2)}{\pi(\pi^2+3)}. &
\end{flalign*}

\section*{7. Множители и экстремаль}

Найденные множители Лагранжа:
\begin{flalign*}
& \lambda_0=1,\qquad \lambda_1=-4a=-\frac8{\pi^2+3},\qquad \lambda_2=4b=\frac{8(\pi^2+4)}{\pi(\pi^2+3)}, & \\
& \lambda_3=p(0)=2\dot x(0)=-\frac{4\pi}{\pi^2+3}. &
\end{flalign*}

Единственная допустимая экстремаль:
\begin{flalign*}
& \hat x(t)=\frac2{\pi(\pi^2+3)}\Bigl[-2(\pi^2+2)\sin t+\pi\,t\sin t+(\pi^2+4)\,t\cos t\Bigr]. &
\end{flalign*}

Значение функционала на экстремали:
\begin{flalign*}
& J(\hat x)=-\pi a-4b=-\frac{2(5\pi^2+16)}{\pi(\pi^2+3)}\approx-3{,}2326. &
\end{flalign*}

\section*{8. Исследование достаточности}

Квадратичный функционал второй вариации (приращение):
\begin{flalign*}
& \Delta J(h) = \int_0^\pi \bigl(\dot{h}^2 - h^2\bigr) \, dt. &
\end{flalign*}

\begin{enumerate}
    \item \textbf{Условие Лежандра:}
    \begin{flalign*}
    & F_{\dot x\dot x} = 2\lambda_0 = 2 > 0 \implies \text{кандидат на строгий локальный минимум}. &
    \end{flalign*}

    \item \textbf{Условие Якоби.}
    Уравнение Якоби:
    \begin{flalign*}
    & \ddot y + y = 0, \quad y(0)=0, \quad \dot y(0)=1 \implies y(t) = \sin t. &
    \end{flalign*}
    \begin{flalign*}
    & y(\pi) = 0 \implies t^* = \pi \text{ — сопряжённая точка.} &
    \end{flalign*}
    \begin{itemize}
        \item Обычное условие Якоби (на интервале \((0,\pi)\)): \textbf{выполнено}.
        \item Усиленное условие Якоби (на полуинтервале \((0,\pi]\)): \textbf{не выполнено}.
    \end{itemize}
    
    Без ограничений: \(h = \sin t \implies \Delta J(\sin t) = 0\).
    
    Для допустимых вариаций:
    \begin{flalign*}
    & \int_0^\pi h\cos t \, dt = 0, \quad \int_0^\pi h\sin t \, dt = 0 \implies \sin t \notin \{h\} \implies \Delta J(h) > 0 \quad (h \not\equiv 0). &
    \end{flalign*}

    \item \textbf{Проверка вариацией.} Возьмем \(h = \sin 3t\):
    \begin{flalign*}
    & \int_0^\pi \sin 3t\cos t \, dt = 0, \quad \int_0^\pi \sin 3t\sin t \, dt = 0 \quad (\text{допустима}). &
    \end{flalign*}
    \begin{flalign*}
    & \Delta J(\sin 3t) = \int_0^\pi \bigl(9\cos^2 3t - \sin^2 3t\bigr) \, dt = 4\pi > 0. &
    \end{flalign*}
    \begin{flalign*}
    & J(\hat x + \mu \sin 3t) = J(\hat x) + 4\pi\mu^2 \xrightarrow{\mu\to\infty} +\infty \implies \text{глобального максимума нет.} &
    \end{flalign*}
\end{enumerate}

\textbf{Вывод:} \(\Delta J(h) > 0\) для любых допустимых \(h \not\equiv 0 \implies \hat x(t)\) доставляет \textbf{строгий глобальный минимум}.

\vspace{1cm}
\section*{Ответ}

\begin{flalign*}
& \hat x(t)=\frac2{\pi(\pi^2+3)}\Bigl[-2(\pi^2+2)\sin t+\pi\,t\sin t+(\pi^2+4)\,t\cos t\Bigr], & \\
\\& J_{\min}=-\frac{2(5\pi^2+16)}{\pi(\pi^2+3)}\approx-3{,}2326. &
\end{flalign*}

Траектория \(\hat x(t)\) доставляет \textbf{строгий глобальный минимум}; глобального максимума не существует.

\end{document}
